% 图 1
\documentclass[tikz,border=4pt]{standalone}
\usepackage{amsmath}
\usetikzlibrary{arrows.meta,calc}

\definecolor{memberred}{HTML}{B20000}

\begin{document}
\begin{tikzpicture}[
  line cap=round,
  line join=round,
  force/.style={
    draw=black,
    line width=0.9pt,
    -{Stealth[length=2.8mm,width=2.2mm]}
  },
  guide/.style={
    draw=black,
    line width=0.8pt,
    dash pattern=on 2.2pt off 2.2pt
  },
  anglemark/.style={
    draw=black,
    line width=0.8pt
  },
  every node/.style={font=\small}
]

  \coordinate (A) at (0,0);
  \coordinate (C) at (1.4,1.4);
  \coordinate (B) at (2.8,2.8);

  % 红色连杆
  \draw[memberred,line width=1.0pt] (A) -- (B);

  % 上端节点 i
  \draw[guide] (B) -- ++(0,1.40);
  \draw[guide] (B) -- ++(0,-0.70);
  \draw[guide] (B) -- ++(1.40,1.40);

  \draw[force]
    (B) -- ++(1.05,2.10)
    node[above right=-1pt] {$T_i$};

  % 上端角度
  \draw[anglemark]
    ($(B)+(90:1.05)$)
    arc[
      start angle=90,
      end angle=63.435,
      radius=1.05
    ];

  \draw[anglemark]
    ($(B)+(90:0.875)$)
    arc[
      start angle=90,
      end angle=45,
      radius=0.875
    ];

  \node at ($(B)+(0.35,1.23)$) {$\theta_i$};
  \node at ($(B)+(0.70,0.53)$) {$\varphi_i$};

  % 中部受力
  \draw[force]
    (C) -- ++(0,1.05)
    node[above] {$F_i$};

  \draw[force]
    (C) -- ++(0,-1.40)
    node[below] {$G_i$};

  \draw[anglemark]
    ($(C)+(-90:0.35)$)
    arc[
      start angle=-90,
      end angle=-135,
      radius=0.35
    ];

  \node at ($(C)+(-0.18,-0.53)$) {$\varphi_1$};

  % 下端节点 i+1
  \draw[guide]
    ($(A)+(0,0.70)$) -- ($(A)+(0,-1.40)$);

  \draw[guide]
    (A) -- ++(-1.40,-1.40);

  \draw[force]
    (A) -- ++(-2.45,-0.70)
    node[below left=-1pt] {$T_{i+1}$};

  % 下端角度
  \draw[anglemark]
    ($(A)+(-90:1.05)$)
    arc[
      start angle=-90,
      end angle=-135,
      radius=1.05
    ];

  \draw[anglemark]
    ($(A)+(-90:0.70)$)
    arc[
      start angle=-90,
      end angle=-164.055,
      radius=0.70
    ];

  \node at ($(A)+(-0.88,-0.53)$) {$\theta_{i+1}$};
  \node at ($(A)+(-0.53,-1.23)$) {$\varphi_i$};

  \path[use as bounding box]
    (-3.25,-1.75) rectangle (4.65,5.95);

\end{tikzpicture}
\end{document}